\section{Diskussion}\label{sec:discussion}

Diskussionskapitlet inkluderar en diskussion kring resultaten som uppnåts i projektet, samt en analys av metoder, tekniska förutsättningar och vidare utveckling av systemen.

\subsection{Analys av Resultat}\label{subsec:discussion_analysis}

Projektets huvudsakliga mål var att undersöka och utveckla en lösning som med hjälp av AI effektiviserar informationssökning inom en organisation.
Baserat på resultaten kan det konstateras att bägge de testade lösningarna fungerar på olika plan.

Ett tekniskt hinder var att mätningarna inte gjordes på samma mängd data, samt data med olika innehåll.

\subsubsection{Graf-baserade lösningen}

Den graf-baserade lösningen presterade överlag väl och visade tydliga fördelar i de scenarion den var designad för.
Resultaten visar ett mönster där lösningens styrkor framträder mer och mer när frågorna blir komplexare och data från flera källor behöver kombineras.

I de enklare frågorna behövdes det generellt sett göras mer verktygsanrop än den api-baserade lösningen. Det tog även längre tid från en fråga ställs till att ett svar är genererat.
Anledningen till att det tar längre tid är för att det blir en del overhead då LLM:en först måste hämta hela graf-schemat och sedan generera en cypher fråga som exekveras mot databasen.
Det blir fler verktygsanrop än API-baserade lösningen då min design av lösningen alltid kallar på verktyget som returnerar schemat för grafen, även för de mest simpla frågorna.

I de medelsvåra frågorna där det behövdes mer resonemang presterade lösningen bra. Den svarade rätt på alla frågor förutom en där den egentligen hade delvis rätt.
Frågan den hade fel på var fråga nio där det frågades om det fanns några blockers i projektet. Med blockers i projektet menade jag egentligen bara vanliga issues.
Men LLM:en resonerade som att en blocker var en issue med label \texttt{blocked}, \texttt{stuck}, \texttt{at risk} eller \texttt{risk}, se Figur \ref{fig:kg_q9}.
Där tog jag beslutet att markera svaret som felaktigt eftersom jag räknar alla issues som blockers.

I de komplexa frågorna presterade lösningen bra. Den hade väldigt få verktygsanrop (lika många som i de simpla frågorna) och svarade rätt på alla frågor förutom en.
Tiden det tog att svara på på frågorna var ungefär lika lång som när den svarade på de simpla frågorna.
I dessa frågor behövdes det resoneras kring hur data hänger ihop mellan olika källor, och det är just det graflösningen är designad för att göra.
Genom att det finns kopplingar mellan exempelvis en person till ett repository och från samma person till ett slack inlägg kan LLM:en enkelt följa dessa kopplingar för att se att det är samma person som har gjort ett inlägg som är med i ett repository.

Ett bevis på detta resonemang ser vi i fråga 15 där det frågades om namnet på det senast uppdaterade confluence dokumentet från personen som gjorde den senaste commiten i repositoryt "payment-service".
Som vi ser i Figur \ref{fig:kg_q15}, frågar LLM:en först efter schemat, sedan efter personen som gjorde den senaste commiten i det efterfrågade repositoryt, och till sist efter det senaste uppdaterade dokumentet från den personen.
Den hämtade denna informationen med en enda generard cypher fråga, om man räknar bort schema anropet.

Det är i frågor som denna vi ser styrkan i den graf-baserade lösningen, då den ger LLM:en den kontext den behöver för att kunna följa samband mellan data från olika källor.

\subsubsection{API-baserade Lösning}

Den API-baserade lösningen presterade bra på de enklare frågorna där lika mycket resonemang inte behövdes, men visade tydliga svagheter i de mer komplexa frågorna där kopplingar mellan olika datakällor behövde göras.

I de enklare frågorna presterade lösningen bra och svarade rätt på alla frågor med i genomsnitt färre verktygsanrop och med lite kortare tid än den graf-baserade lösningen.
Det var förväntat att den här sortens lösning skulle prestera bra när det är tydliga frågor som endast kräver data från en källa.
 Verktygen är designade för att hämta specifik information från en källa och i de enklare frågorna var det endast det som behövdes.

I de medelsvåra frågorna började lösningen visa svagheter. Den svarade rätt på 40\% av frågorna trots att den i genomsnitt hade fler verktygsanrop.
Det var lite överraskande att den hade så många fel i denna komplexitet då dessa frågor inte handlar om att koppla ihop data från specifika källor utan mer att hämta olika data från samma källa och sedan resonera kring det.

Att den presterade så dåligt i denna kategori beror troligen på att frågorna krävde tolkningar och filtreringar som verktygen inte stödjer.
LLM:en fick råa svar från API:erna och behövde sedan själv resonera kring vad som var relevant, vilket då i flera fall ledde till felaktiga slutsatser.

I de komplexa frågorna presterade lösningen dåligt med endast 20\% korrekta svar.
Lösningen hade stora problem med att koppla data från en källa med data från en annan.
Till skillnad från den graf-baserade lösningen där en person representeras som en nod med tydliga relationer till andra entiteter, kan samma person ha olika namn eller representation i olika system i den API-baserade lösningen, vilket gör det svårt för LLM:en att avgöra att det är samma person.
Detta är en bidragande orsak till att lösningen svarade fel på det mesta.

Detta bekräftar hypotesen om att en relationsbaserad graf med tydliga samband mellan entiteter underlättar för LLM:en att navigera och hämta data från olika källor.

\subsection{Metoddiskussion}\label{subsec:project_method_discussion}
Make a deep analysis and discussion of your chosen method, chosen approach,
chosen metrics, etc. Answer all of your project phases/milestones. Have they
been met?

\subsection{Vetenskaplig diskussion}\label{subsec:scientific_discussion}
Make a deep discussion of the gained scientific knowledge and your scientific
method. What can be learnt? Is this knowledge general or specific, etc. Make a
discussion by looking back at the related works and put your work into
perspective. Remember to discuss and show insights into the scientific
possibilities of this work and its limitations

\subsection{Rekommendationer och konsekvenser}\label{subsec:recommendation} 

Discuss the scientific impact and the contribution of your work. What will be the consequences of
this new knowledge? If you are working with an industry problem, what do you recommend the company
to continue to pursue?

\subsection{Etiska och samhälleliga aspekter}\label{subsec:ethical_discussions}
You will need to include a discussion on ethics, societal impact, and
considerations. Use the human perspective, how will we be effected by this work,
was people involved in the process, privacy? Remember to discuss and show
insights into this projects role in society and our responsibilities for how it
is used.

\subsection{Verklig distribution av system}\label{subsec:real_system_distribution}

